\section{Related Work}
\label{sec:related}

\paragraph{Manually designed kernels.}
Existing ML frameworks such as TensorFlow XLA~\cite{tensorflow_xla, Tensorflow}, PyTorch~\cite{pytorch}, and TensorRT~\cite{tensorrt} adopt a kernel-per-operator approach and rely on GPU experts to manually design and implement kernels for individual operators.
For attention alone, various specialized kernels have been developed, including FlashAttention~\cite{dao2023flash, tri2023flashdecoding, hong2024flashdecoding}, FasterTransformer~\cite{fastertransformer}, and FlashInfer~\cite{ye2025flashinfer}, each targeting specific architectural features or usage scenarios. Current systems rely on a fragmented ecosystem of specialized kernel libraries, making it difficult to unify the entire inference pipeline within a single, holistic mega-kernel.
% cohesive mega-kernel.

%\paragraph{Manually-designed kernels.} Existing ML frameworks such as TensorFlow XLA~\cite{tensorflow_xla, Tensorflow}, PyTorch~\cite{pytorch}, and TensorRT~\cite{tensorrt} takes a kernel-per-operator approach and relies on GPU experts to manually design kernels for individual ML operators.
%
%For example, to accelerate attention~\cite{transformers}, several specialized kernels have been developed based on FlashAttention~\cite{dao2023flash, tri2023flashdecoding, hong2024flashdecoding, fastertransformer, ye2025flashinfer}. 
%Current frameworks reply on a fragmented ecosystem of specialized kernel libraries, making it difficulty to unify the entire inference pipeline with in a single mega-kernel.

\paragraph{ML compilers.}
A large body of work has explored {\em compiler-based} generation of high-performance kernels for tensor programs. Systems such as TVM~\cite{tvm, tvm_auto_tuner}, Ansor~\cite{ansor}, and Triton~\cite{tillet2019triton}, alongside others~\cite{zheng2020flextensor, hagedorn2023graphene, feng2022tensorir, hu2024korch}, build on the algorithm–schedule separation introduced in Halide~\cite{halide, autohalide}.
Another line of work employs {\em superoptimization} techniques to automatically search for efficient kernel implementations from high-level specifications~\cite{TASO, wang2021pet, zheng2023einnet, tensat, unger2022unity}.
However, these compilers are fundamentally designed around operator-level optimization and do not support generating a unified mega-kernel or coordinating cross-operator execution.

%\paragraph{ML compilers.} Recent work has introduced a variety of ML compilers that generate high-performance kernels for tensor programs. 
%TVM~\cite{tvm, tvm_auto_tuner}, Ansor~\cite{ansor}, and Triton~\cite{tillet2019triton}, along with others~\cite{zheng2020flextensor, hagedorn2023graphene, feng2022tensorir, hu2024korch} leverages the idea of algorithm-schedule separation introduced in Halide~\cite{halide, autohalide}.
%Recent work has also applied superoptimization techniques to automatically search for fast implementation~\cite{TASO, wang2021pet, zheng2023einnet, tensat, unger2022unity}.
%These compilers are designed and optimized for individual operators and do not support mega-kernel generation.

\paragraph{Mega-kernels.}
Prior efforts on mega-kernels largely rely on manual design. For example, FlashDMoE fuses mixture-of-experts computation with inter-GPU communication into a single handcrafted mega-kernel~\cite{aimuyo2025flashdmoe}. As another example, Spector et al.\ manually developed a low-latency mega-kernel for LLaMA-1B~\cite{touvron2023llama, megakernel}.
These approaches require extensive engineering effort and deep GPU expertise, and they do not generalize across models or GPUs.
In contrast, \sys provides a compiler-based solution that automatically transforms a tensor program into a highly optimized mega-kernel, eliminating the need for manual fusion or hand-written mega-kernel implementations.

%\paragraph{Mega-kernels.}
%Existing mega-kernels are handcrafted by GPU experts for specific models.
%
%For example, FlashDMoE fuses mixture-of-experts computation and inter-GPU communication into a single kernel~\cite{aimuyo2025flashdmoe}, while Spector et al. manually designed and implemented a low-latency mega-kernel for LLAMA-1B~\cite{touvron2023llama}. 
%These manual approaches require substantial engineering effort and deep GPU expertise to mega-kernelize a tensor program. In contrast, \sys adopts a compiler-based approach that automatically transforms an input tensor program into an optimized mega-kernel, eliminating the need for manual design or implementation. 
